%%%%%%   TIPO DE DOCUMEN%TO: Reporte   %%%%%%
\documentclass[letterpaper,11pt]{report}
%%%%%%%%%%%%%%%%%%%%%%%%%%%%%%%%%%%%%%%%%%%%

%%%%%%%   PAQUETES   %%%%%%%
\usepackage[spanish]{babel}
\usepackage{graphicx}
\usepackage[utf8]{inputenc}
%%%%%%%%%%%%%%%%%%%%%%%%%%%%

%%%%%%%%%%%%%%%%%%%%%%%%%%%%%%%%%%%%%%%%
%%%%%%%   INICIO DEL DOCUMENTO   %%%%%%%
%%%%%%%%%%%%%%%%%%%%%%%%%%%%%%%%%%%%%%%%
\begin{document}

    %%%%%%% Renombrar en espaNol %%%%%%
    \renewcommand{\tablename}{Tabla} %Escribe Tabla en lugar de Cuadro
    \renewcommand{\listtablename}{\'Indice de tablas} %Escribe Indeice de tablas en lugar de Indice de cuadros

    %%%%%%%   PORTADA   y RESUMEN   %%%%%%% Este es el contenido agregado de la secci\'on 2.

    %%%%%%%%%%%%%%%%%%%%%%%%%%%%%
%%%%%      PORTADA      %%%%%
%%%%%%%%%%%%%%%%%%%%%%%%%%%%%

\begin{titlepage}

    \centering %Todo centrado

    %%%%  LOGO DE LA ESCUELA   %%%%
    \includegraphics[scale=0.17]{imagenes/escom-ipn} %Imagen para portada
    %%%%  NOMBRE DE LA ESCUELA   %%%%
    \LARGE{\\ Instituto Polit\'ecnico Nacional}
    \LARGE{\\ Escuela Superior de C\'omputo}
    
    \vspace{1cm} %Espacio vertical

    %%%%  TITULO Y NÚMERO DE TRABAJO   %%%%
    \LARGE \textbf{ Nombre del Trabajo Terminal}
    \LARGE {\\ Número de Trabajo Terminal}

    \vspace{1cm} %Espacio vertical

    \LARGE \textit{Que para cumplir con la opción de titulación curricular en la carrera de:}
    \LARGE \textbf{\\ Ingeniería en Sistemas Computacionales}

    \vspace{1cm} %Espacio vertical

    %%%%   ALUMNOS   %%%%
   \textit{Presentan}\\
    Cesar Raúl Avila Padilla \\
    Ivo Sebatians Sam-Álvarez Tostado\\

    \vspace{1cm} %Espacio vertical

    %%%%   Directores   %%%%
   \textit{Directores}\\
   M en C. Ulises Velez Saldaña  \bigskipM. en C. José David Ortega Pacheco 
   

\end{titlepage} %incluye el archivo portada.tex
    %%%%%%%%%%%%%%%%%%%%%%%
%%%%    RESUMEN    %%%%
%%%%%%%%%%%%%%%%%%%%%%%

\begin{abstract}

  Aquí va el resumen general del documento de trabajo terminal.

\end{abstract}
 %Incluir resumen del documento (resumen.tex)

    %%%%%%%   INCLUIR ENCABEZADOS EN INDICES Y CAPITULOS   %%%%%%%
    \pagestyle{headings}

    %%%%%%%   NUMERACION EN CONTENIDO E INDICE DE TABLAS Y FIGURAS   %%%%%%%
    \pagenumbering{roman} %NUmeros romanos
    %\setcounter{page}{1} %Comienza en I por default, aquI se puedo modificar

    %%%%%%   INCLUIR CONTENIDO, INDICE DE FIGURAS E INDICE DE TABLAS   %%%%%%
    \tableofcontents
    \listoffigures
    \listoftables

    %%%%%%%   NUMERACION EN CAPITULOS   %%%%%%%
    \clearpage %Para iniciar con los arAbigos
    \pagenumbering{arabic} %Numeros arabigos
    %\setcounter{page}{1} %Comienza en 1 por default, aquI se puede modificar

    %%%%%%%   INCLUYE CAPITULOS Y SECCIONES   %%%%%%%
    \chapter{Introducci\'on}

Colocar una descripción de lo que contiene el capítulo
	   \section{Contexto de trabajo}

Se describe el área donde se esta trabajando y en qué contexto
	   \section{Problemática}

   
	   \section{Trabajo previo}

  
	   \section{Solución propuesta}

 
	   \section{Objetivos}

	   \section{Justificación}


    \chapter{Tema Teórico 1} %CAP\'ITULO 2

	   \section{Subtema 1}


	   \section{Subtema 2}


%    \input{./temaTeorico2/temaTeorico2}
%	   \section{Subtema 1}


%	   \section{Subtema 2}





%    \input{./marcoTeorico/marco} %CapItulo 2 (capitulo2.tex)
%	   \input{./marcoTeorico/secciones/eyeTracking}

%    \input{./modeloDeComportamiento/modeloDeComportamiento}
%    \input{./modeloDeComportamiento/requerimientosGenerales}
%    \input{./modeloDeComportamiento/modulos}

%    \input{./modeloDeComportamiento/modeloDeComportamiento}



%    \input{./modeloInteraccion/modeloInteraccion}
%    \input{./modeloInteraccion/DisenoPantallas}
%    \input{./modeloInteraccion/catalogoMensajes}

\end{document}
%%%%%%%%%%%%%%%%%%%%%%%%%%%%%%%%%%%%%%%
%%%%%%%    FIN DEL DOCUMENTO    %%%%%%%
%%%%%%%%%%%%%%%%%%%%%%%%%%%%%%%%%%%%%%%
